\chapter{Documentos técnicos e escrita matemática}

Depois de entender o ciclo de compilação, os arquivos auxiliares e o papel dos pacotes, o próximo passo é usar essa infraestrutura para escrever documentos técnicos que continuam legíveis quando crescem. Um texto matemático pequeno tolera improvisos: uma equação isolada, um comando escrito diretamente no corpo do texto, uma referência feita à mão. Em um artigo, uma dissertação, uma lista extensa ou uma apostila, esses improvisos passam a gerar inconsistências.

Este capítulo trata de três problemas recorrentes: organizar o documento em partes, compor matemática com os ambientes apropriados e criar comandos semânticos para a notação usada no texto. O objetivo não é decorar todos os comandos dos pacotes da American Mathematical Society, mas entender quais decisões eles permitem tomar.

\section{Estrutura de um documento técnico}

Um documento técnico tende a crescer em camadas. O preâmbulo reúne decisões globais: classe, pacotes, fontes, idioma, estilo de referências e comandos próprios. O corpo do documento contém o texto. Arquivos auxiliares são gerados pelo compilador e registram informações necessárias entre passagens, como labels, citações e sumário.

Misturar todas essas camadas no mesmo arquivo funciona em exemplos mínimos, mas dificulta manutenção. A dificuldade aparece quando o autor precisa mover uma seção, revisar apenas um capítulo, localizar a definição de uma macro ou comparar versões do texto. Um documento técnico organizado deve deixar claro onde ficam as decisões globais e onde fica o conteúdo.

Uma estrutura simples e suficiente para muitos projetos é:

\begin{lstlisting}[caption={Estrutura possível de projeto}]
artigo/
  artigo.tex
  secoes/
    introducao.tex
    preliminares.tex
    resultados.tex
  figs/
    dominio.pdf
    grafico.pdf
  macros.tex
  referencias.bib
\end{lstlisting}

Essa organização não é obrigatória, mas explicita responsabilidades. O arquivo principal controla a compilação. Os arquivos de seção contêm texto. Figuras ficam em diretório próprio. Macros recorrentes podem ficar no preâmbulo ou em um arquivo dedicado, carregado antes do corpo do documento.

Uma forma simples de modularização é manter um arquivo principal e carregar capítulos ou seções com \verb|\input|:

\begin{lstlisting}[language={[latex]TeX}, caption={Arquivo principal com capítulos separados}]
\documentclass{book}

\usepackage[T1]{fontenc}
\usepackage[utf8]{inputenc}
\usepackage[brazilian]{babel}
\usepackage{mathtools}
\usepackage{amssymb}

\begin{document}

\input{introducao}
\input{preliminares}
\input{resultados}

\end{document}
\end{lstlisting}

O comando \verb|\input| copia o conteúdo do arquivo indicado naquele ponto da compilação. Para o leitor e para o compilador, o resultado é como se o conteúdo estivesse escrito no arquivo principal. A vantagem é organizacional: cada capítulo pode ser editado separadamente, e o arquivo principal continua pequeno o bastante para mostrar a estrutura geral do projeto.

O comando \verb|\include| também carrega arquivos, mas cria quebras de página e trabalha com arquivos auxiliares próprios. Ele é útil em livros ou teses grandes, principalmente quando combinado com \verb|\includeonly| para compilar apenas partes do projeto durante a revisão.

\begin{lstlisting}[language={[latex]TeX}, caption={Compilação parcial com includeonly}]
\includeonly{capitulos/resultados}

\begin{document}
\include{capitulos/introducao}
\include{capitulos/preliminares}
\include{capitulos/resultados}
\end{document}
\end{lstlisting}

Em documentos técnicos, nomes previsíveis ajudam. Arquivos como \texttt{preliminares.tex}, \texttt{resultados.tex} e \texttt{apendice-demonstracoes.tex} comunicam melhor o papel do arquivo do que nomes genéricos como \texttt{parte1.tex}. O mesmo vale para labels: \verb|sec:preliminares|, \verb|eq:energia| e \verb|fig:dominio| são mais úteis do que \verb|label1|. Uma convenção pequena, aplicada de forma consistente, reduz o custo de manutenção.

\section{A pilha AMS}

O pacote \texttt{amsmath} é a base moderna para escrita matemática em \LaTeX. Ele fornece ambientes de equações, comandos para operadores, controle de numeração, matrizes, casos, setas extensíveis, frações e vários ajustes tipográficos. O pacote \texttt{amssymb} adiciona símbolos matemáticos comuns em textos avançados. O pacote \texttt{mathtools} estende o \texttt{amsmath}: carrega \texttt{amsmath} automaticamente e oferece ferramentas para delimitadores, versões refinadas de ambientes e comandos úteis para ajustes finos.

É útil pensar nesses pacotes como camadas:

\begin{description}
    \item[\texttt{amsmath}] fornece a infraestrutura principal para fórmulas exibidas, alinhamentos, numeração, operadores e construções matemáticas.
    \item[\texttt{amssymb}] amplia o repertório de símbolos, como relações, setas e alfabetos matemáticos.
    \item[\texttt{mathtools}] complementa o \texttt{amsmath} com recursos práticos para delimitadores, casos em estilo de display, interrupções curtas em alinhamentos e ajustes finos de espaçamento.
\end{description}

Uma configuração comum para documentos com bastante matemática é:

\begin{lstlisting}[language={[latex]TeX}, caption={Pacotes para escrita matemática}]
\usepackage{mathtools} % carrega amsmath
\usepackage{amssymb}
\end{lstlisting}

Também é comum encontrar:

\begin{lstlisting}[language={[latex]TeX}, caption={Carregamento explícito de amsmath}]
\usepackage{amsmath}
\usepackage{amssymb}
\end{lstlisting}

As duas escolhas são aceitáveis. Quando o documento usa recursos de \texttt{mathtools}, carregar apenas \verb|\usepackage{mathtools}| evita redundância. Quando o documento usa apenas recursos clássicos da AMS, \verb|\usepackage{amsmath}| deixa a dependência explícita.

Em um minicurso avançado, a questão mais importante não é ``qual pacote carrega mais comandos'', mas qual pacote permite escrever a intenção matemática com menos improviso. Ambientes como \texttt{align} e \texttt{split} descrevem a estrutura da dedução. Comandos como \verb|\DeclareMathOperator| descrevem operadores. \verb|\DeclarePairedDelimiter| descreve delimitadores que têm papel semântico no texto.

\section{Ambientes de equações}

Um erro comum é tratar todos os displays matemáticos como se fossem iguais. A escolha do ambiente comunica a estrutura da fórmula: uma equação isolada, uma equação longa quebrada em linhas, um grupo de equações independentes ou um cálculo alinhado por relações.

Para uma equação única numerada, use \texttt{equation}. Para uma equação única sem número, use \verb|\[...\]| ou \texttt{equation*}:

\begin{lstlisting}[language={[latex]TeX}, caption={Equação isolada}]
\begin{equation}\label{eq:energia}
E = mc^2
\end{equation}

\[
E = mc^2
\]
\end{lstlisting}

Quando uma equação é longa, mas ainda representa uma única igualdade ou identidade, \texttt{split} dentro de \texttt{equation} permite quebrar linhas mantendo um único número:

\begin{lstlisting}[language={[latex]TeX}, caption={Equação longa com split}]
\begin{equation}\label{eq:funcional}
\begin{split}
J(u)
  &= \int_\Omega |\nabla u(x)|^2\,dx \\
  &\quad + \int_\Omega V(x)|u(x)|^2\,dx
       - \int_\Omega F(x,u(x))\,dx .
\end{split}
\end{equation}
\end{lstlisting}

O ponto de alinhamento vem antes do símbolo relacional. Em \texttt{align}, por exemplo, escreve-se \verb|&=|, não \verb|=&|. Isso preserva o espaçamento matemático esperado ao redor da relação.

\begin{lstlisting}[language={[latex]TeX}, caption={Cálculo alinhado}]
\begin{align}
\|u+v\|^2
  &= \langle u+v, u+v\rangle \\
  &= \|u\|^2 + 2\langle u,v\rangle + \|v\|^2 .
\end{align}
\end{lstlisting}

Quando há várias equações sem alinhamento comum, \texttt{gather} é mais adequado que \texttt{align}. Quando há uma equação longa sem ponto natural de alinhamento, \texttt{multline} pode ser melhor que forçar um alinhamento artificial.

\begin{lstlisting}[language={[latex]TeX}, caption={Equações sem alinhamento comum}]
\begin{gather}
A+B+C = 0,\\
\int_\Omega f(x)\,dx = 1.
\end{gather}
\end{lstlisting}

Para equações longas sem uma coluna de alinhamento natural, \texttt{multline} distribui a primeira linha à esquerda, a última à direita e as linhas intermediárias de forma centrada:

\begin{lstlisting}[language={[latex]TeX}, caption={Equação longa sem alinhamento natural}]
\begin{multline}
a_1+a_2+\cdots+a_n
+ b_1+b_2+\cdots+b_n \\
+ c_1+c_2+\cdots+c_n
+ d_1+d_2+\cdots+d_n = 0.
\end{multline}
\end{lstlisting}

Quando o espaçamento entre colunas precisa ser controlado explicitamente, \texttt{alignat} é uma alternativa a \texttt{align}. Ele exige informar quantos pares de colunas alinhadas serão usados:

\begin{lstlisting}[language={[latex]TeX}, caption={Alinhamento com controle de espaçamento}]
\begin{alignat}{2}
u_t &= \Delta u &&+ f(u),\\
v_t &= \Delta v &&+ g(v).
\end{alignat}
\end{lstlisting}

O ambiente \texttt{eqnarray} deve ser evitado em documentos novos. Ele produz espaçamento inadequado ao redor de relações e lida mal com numeração em equações longas. Para documentos já existentes, a substituição por \texttt{align} ou por \texttt{equation} com \texttt{split} costuma ser uma melhoria objetiva.

\section{Numeração e referências}

Referências para equações devem usar labels, não números digitados manualmente. O comando \verb|\eqref|, fornecido por \texttt{amsmath}, imprime a referência com parênteses:

\begin{lstlisting}[language={[latex]TeX}, caption={Referência a equação}]
\begin{equation}\label{eq:identidade}
(a+b)^2 = a^2 + 2ab + b^2
\end{equation}

Pela identidade \eqref{eq:identidade}, obtemos...
\end{lstlisting}

Em ambientes com várias linhas, nem toda linha precisa ser numerada. Use \verb|\notag| quando a linha não deve receber número:

\begin{lstlisting}[language={[latex]TeX}, caption={Controle de numeração por linha}]
\begin{align}
f(x) &= x^2 + 2x + 1 \notag\\
     &= (x+1)^2. \label{eq:quadrado}
\end{align}
\end{lstlisting}

Quando um grupo de equações deve compartilhar uma numeração principal, com letras subordinadas, use \texttt{subequations}:

\begin{lstlisting}[language={[latex]TeX}, caption={Numeração subordinada}]
\begin{subequations}\label{eq:sistema}
\begin{align}
-\Delta u &= f && \text{em } \Omega,\\
u &= 0 && \text{em } \partial\Omega.
\end{align}
\end{subequations}
\end{lstlisting}

Esse padrão é útil quando o texto precisa referir-se ao sistema inteiro e também a cada equação individual.

Em textos longos, pode ser desejável numerar equações por seção ou capítulo. O \texttt{amsmath} fornece \verb|\numberwithin| para amarrar o contador de equações a outro contador:

\begin{lstlisting}[language={[latex]TeX}, caption={Numeração por seção}]
\numberwithin{equation}{section}
\end{lstlisting}

Com isso, equações podem aparecer como (2.1), (2.2), (3.1) e assim por diante. Essa escolha deve ser feita no preâmbulo e precisa ser coerente com o tipo de documento. Em um artigo curto, numeração sequencial pode ser suficiente. Em uma tese, livro ou apostila extensa, numeração por seção ou capítulo facilita a localização.

Há ainda situações em que um número automático não é a melhor etiqueta. O comando \verb|\tag| permite indicar uma marca manual, por exemplo em uma hipótese ou identidade clássica:

\begin{lstlisting}[language={[latex]TeX}, caption={Etiqueta manual com tag}]
\begin{equation}\tag{H}\label{eq:hipotese-H}
\int_\Omega F(x,u)\,dx \leq C(1+\|u\|^p).
\end{equation}
\end{lstlisting}

Esse recurso deve ser usado com moderação. Se tudo recebe uma etiqueta manual, o documento perde o benefício da numeração automática.

\section{Texto dentro de matemática}

Texto dentro de uma fórmula não deve ser improvisado com \verb|\mathrm| quando a intenção é escrever uma palavra ou uma frase. O comando \verb|\text| preserva o tamanho adequado em índices e mantém o texto como texto:

\begin{lstlisting}[language={[latex]TeX}, caption={Texto em uma fórmula}]
\[
P_{r-j} =
\begin{cases}
0, & \text{se } r-j \text{ e impar},\\
1, & \text{se } r-j \text{ e par}.
\end{cases}
\]
\end{lstlisting}

Em cálculos alinhados, justificativas curtas podem aparecer em uma coluna separada:

\begin{lstlisting}[language={[latex]TeX}, caption={Justificativas em align}]
\begin{align}
x &= y_1 - y_2 + y_3 && \text{por \eqref{eq:recorrencia}}\\
  &= y^\prime \circ y^* && \text{pela definicao de composicao}.
\end{align}
\end{lstlisting}

Quando uma frase precisa interromper um cálculo mantendo o alinhamento, \verb|\intertext| é apropriado. Para uma interrupção muito curta, \texttt{mathtools} oferece \verb|\shortintertext|, que usa menos espaço vertical.

\begin{lstlisting}[language={[latex]TeX}, caption={Interrupção preservando alinhamento}]
\begin{align}
A_1 &= N_0(\lambda;\Omega') - \phi(\lambda;\Omega'),\\
A_2 &= \phi(\lambda;\Omega') - \phi(\lambda;\Omega),
\shortintertext{e portanto}
A_3 &= N(\lambda;\omega).
\end{align}
\end{lstlisting}

\section{Operadores matemáticos}

Funções como \verb|\sin|, \verb|\log| e \verb|\lim| são compostas em fonte romana e recebem espaçamento próprio. O mesmo tratamento deve ser dado a operadores definidos pelo autor. Escrever \verb|rank A|, \verb|\mathrm{rank} A| ou \verb|Rank A| diretamente no corpo do documento resolve apenas o caso local; declarar o operador resolve a notação do documento inteiro.

\begin{lstlisting}[language={[latex]TeX}, caption={Declaração de operadores}]
\DeclareMathOperator{\rank}{rank}
\DeclareMathOperator{\tr}{tr}
\DeclareMathOperator{\id}{id}

\[
\rank A + \rank B \geq \rank(A+B),
\qquad
\tr(AB) = \tr(BA).
\]
\end{lstlisting}

Quando o operador deve receber limites abaixo e acima em modo display, usa-se a forma estrelada:

\begin{lstlisting}[language={[latex]TeX}, caption={Operador com limites}]
\DeclareMathOperator*{\argmin}{arg\,min}
\DeclareMathOperator*{\esssup}{ess\,sup}

\[
u = \argmin_{v\in V} J(v),
\qquad
\esssup_{x\in\Omega} |f(x)| < \infty .
\]
\end{lstlisting}

Essa prática também torna o código mais semântico. O comando \verb|\argmin| diz o que está sendo escrito; uma sequência manual de letras, espaços e fontes apenas tenta reproduzir uma aparência.

\section{Macros semânticas}

Macros próprias são uma das ferramentas mais importantes em documentos matemáticos. Elas reduzem repetição, evitam inconsistências e separam a notação do modo como ela é desenhada. Se uma notação aparece dezenas de vezes, ela deve ter um nome.

\begin{lstlisting}[language={[latex]TeX}, caption={Macros para notação recorrente}]
\newcommand{\R}{\mathbb{R}}
\newcommand{\N}{\mathbb{N}}
\newcommand{\eps}{\varepsilon}
\newcommand{\dd}{\,d}

\[
\int_{\R^n} f(x)\dd x < \eps.
\]
\end{lstlisting}

O nome da macro deve representar o significado pretendido. Em vez de criar apenas abreviações visuais, como \verb|\bbR|, prefira comandos que revelem o papel da notação no texto quando isso for possível. Em textos longos, comandos como \verb|\norm|, \verb|\abs|, \verb|\inner|, \verb|\Prob| e \verb|\E| ajudam a manter consistência.

\section{Definição de comandos}

Criar comandos próprios é parte normal da escrita de documentos técnicos em \LaTeX. O ponto importante é distinguir comandos que representam notação matemática de comandos que apenas encurtam digitação. Uma macro bem escolhida torna o código-fonte mais legível e permite alterar a notação de forma controlada.

A forma usual para criar comandos em documentos \LaTeX\ é \verb|\newcommand|. Ela verifica se o comando já existe e emite erro se houver conflito. Esse comportamento é desejável: se um nome já foi definido por uma classe, por um pacote ou por outra parte do projeto, sobrescrevê-lo silenciosamente pode produzir efeitos difíceis de diagnosticar.

\begin{lstlisting}[language={[latex]TeX}, caption={Definição segura de comandos}]
\newcommand{\R}{\mathbb{R}}
\newcommand{\eps}{\varepsilon}
\newcommand{\norma}[1]{\lVert#1\rVert}

\[
u \in \R^n,
\qquad
\norma{u} < \eps.
\]
\end{lstlisting}

Comandos com argumentos usam marcadores como \verb|#1|, \verb|#2| e assim por diante. No exemplo, \verb|\norma{u}| substitui \verb|#1| por \verb|u|. Para notações mais estruturadas, vale escolher nomes que descrevam a operação:

\begin{lstlisting}[language={[latex]TeX}, caption={Comandos com argumentos}]
\newcommand{\restricao}[2]{#1\vert_{#2}}
\newcommand{\derivada}[2]{\frac{d #1}{d #2}}

\[
\restricao{f}{A},
\qquad
\derivada{y}{x}.
\]
\end{lstlisting}

Quando a intenção é redefinir um comando existente, use \verb|\renewcommand|. Isso deixa explícito que a redefinição é deliberada. Quando o comando deve ser criado apenas se ainda não existir, use \verb|\providecommand|. Essa última forma é comum em arquivos compartilhados entre projetos, porque permite definir uma macro padrão sem sobrescrever uma definição feita pelo documento principal.

\begin{lstlisting}[language={[latex]TeX}, caption={Redefinição e definição condicional}]
\renewcommand{\epsilon}{\varepsilon}
\providecommand{\dd}{\,d}
\end{lstlisting}

Em código TeX mais antigo ou em implementações internas de pacotes, aparece frequentemente o primitivo \verb|\def|. Ele define comandos de forma mais direta e não faz as verificações de segurança de \verb|\newcommand|. Por isso, em material de curso e em documentos de usuário, \verb|\newcommand| e \verb|\renewcommand| são escolhas melhores. \verb|\def| é útil para entender código de baixo nível, mas não deve ser usado como padrão em documentos técnicos comuns.

\begin{lstlisting}[language={[latex]TeX}, caption={Definição primitiva em TeX}]
% Forma primitiva: comum em codigo de pacote,
% mas menos segura para documentos de usuario.
\def\R{\mathbb{R}}
\end{lstlisting}

\section{Cabeçalhos e rodapés}

Cabeçalhos e rodapés fazem parte da estrutura editorial do documento. Eles ajudam o leitor a se localizar, especialmente em textos longos impressos ou em PDFs consultados fora do editor. Em um livro, é comum colocar o número da página no rodapé e informações como título do livro, autor, capítulo ou seção no cabeçalho. Em um artigo, a escolha pode ser mais discreta ou até inexistente, dependendo da classe e das normas de submissão.

Esse assunto merece atenção porque cabeçalho e rodapé ficam em quase todas as páginas. Um erro de configuração pequeno se repete no documento inteiro. Um excesso de informação também. O objetivo não é decorar comandos do pacote, mas projetar uma política de página: o que aparece em páginas comuns, o que aparece em aberturas de capítulo, onde entra o número da página e quais informações realmente ajudam a navegação.

O pacote mais usado para esse controle é \texttt{fancyhdr}. Ele permite limpar o estilo padrão da página, definir campos de cabeçalho e rodapé e controlar as linhas horizontais. Um exemplo mínimo para documento simples é:

\begin{lstlisting}[language={[latex]TeX}, caption={Rodapé simples com número de página}]
\usepackage{fancyhdr}

\pagestyle{fancy}
\fancyhf{}
\fancyfoot[C]{\thepage}
\renewcommand{\headrulewidth}{0pt}
\renewcommand{\footrulewidth}{0pt}
\end{lstlisting}

Nesse caso, o documento usa apenas o número da página centralizado no rodapé. É uma configuração discreta, adequada para textos curtos, listas e materiais em que o cabeçalho não acrescentaria informação útil.

Em documentos longos, especialmente frente e verso, é comum diferenciar páginas pares e ímpares. O exemplo abaixo coloca a marca do capítulo em páginas pares, a marca da seção em páginas ímpares e o número da página no rodapé:

\begin{lstlisting}[language={[latex]TeX}, caption={Cabeçalho e rodapé com fancyhdr}]
\usepackage{fancyhdr}

\pagestyle{fancy}
\fancyhf{} % limpa cabecalho e rodape

\fancyhead[LE]{\nouppercase{\leftmark}}
\fancyhead[RO]{\nouppercase{\rightmark}}
\fancyfoot[C]{\thepage}

\renewcommand{\headrulewidth}{0.4pt}
\renewcommand{\footrulewidth}{0pt}
\end{lstlisting}

As letras usadas por \texttt{fancyhdr} indicam a posição: \texttt{L}, \texttt{C} e \texttt{R} significam esquerda, centro e direita; \texttt{E} e \texttt{O} significam páginas pares e ímpares. Em documentos frente e verso, essa distinção é relevante porque a margem interna muda de lado. O par \texttt{LE}/\texttt{RO}, por exemplo, costuma indicar posições externas: esquerda em página par e direita em página ímpar.

Os comandos \verb|\leftmark| e \verb|\rightmark| são marcas mantidas pela classe do documento. Em classes como \texttt{book}, \verb|\leftmark| normalmente contém informação de capítulo e \verb|\rightmark| contém informação de seção. O comando \verb|\nouppercase| evita que essas marcas sejam convertidas automaticamente para maiúsculas, comportamento comum em alguns estilos.

\begin{lstlisting}[language={[latex]TeX}, caption={Cabeçalhos externos em documento frente e verso}]
\documentclass[twoside]{book}
\usepackage{fancyhdr}

\pagestyle{fancy}
\fancyhf{}

\fancyhead[LE,RO]{\thepage}
\fancyhead[LO]{\nouppercase{\rightmark}}
\fancyhead[RE]{\nouppercase{\leftmark}}
\end{lstlisting}

Esse desenho põe o número de página no lado externo, onde é mais fácil encontrá-lo em material impresso, e deixa a informação textual no lado interno. A escolha inversa também é possível; o importante é que a regra seja consistente.

Também é possível definir um estilo próprio para páginas especiais, como a primeira página de um capítulo. Isso evita que uma página de abertura receba o mesmo cabeçalho de uma página comum:

\begin{lstlisting}[language={[latex]TeX}, caption={Estilo para páginas especiais}]
\fancypagestyle{plain}{%
  \fancyhf{}
  \fancyfoot[C]{\thepage}
  \renewcommand{\headrulewidth}{0pt}
}
\end{lstlisting}

A razão para redefinir \texttt{plain} é que muitas classes aplicam esse estilo automaticamente em páginas de abertura de capítulo. Se o usuário configura apenas \verb|\pagestyle{fancy}|, pode se surpreender ao ver algumas páginas com comportamento diferente. Em vez de combater isso localmente página por página, é melhor definir o estilo \texttt{plain} de forma explícita.

Outro detalhe técnico frequente é a altura do cabeçalho. Se o conteúdo do cabeçalho não cabe na dimensão reservada, o \LaTeX\ pode emitir um aviso envolvendo \verb|\headheight|. A solução não é ignorar o log, mas reservar altura suficiente:

\begin{lstlisting}[language={[latex]TeX}, caption={Ajuste da altura do cabeçalho}]
\setlength{\headheight}{15pt}
\end{lstlisting}

Esse ajuste deve aparecer no preâmbulo, junto da configuração de página. Se o cabeçalho usa fonte maior, duas linhas ou elementos gráficos, pode ser necessário aumentar esse valor. Em um material técnico sóbrio, porém, cabeçalhos de uma linha costumam ser suficientes.

O pacote também permite criar estilos nomeados. Isso é útil quando o documento tem partes com comportamento diferente: capa sem número, páginas preliminares com numeração romana, corpo principal com cabeçalhos completos e apêndices com uma marca simplificada.

\begin{lstlisting}[language={[latex]TeX}, caption={Estilo de página nomeado}]
\fancypagestyle{semcabecalho}{%
  \fancyhf{}
  \fancyfoot[C]{\thepage}
  \renewcommand{\headrulewidth}{0pt}
}

\thispagestyle{semcabecalho}
\end{lstlisting}

\verb|\thispagestyle| altera apenas a página atual. \verb|\pagestyle| altera o estilo a partir daquele ponto. Essa diferença é importante em capas, páginas de abertura, dedicatórias, sumários e páginas com figuras grandes.

Do ponto de vista editorial, cabeçalho e rodapé não devem carregar informação demais. Se o cabeçalho fica mais saliente que o texto, ele deixa de orientar e passa a competir com o conteúdo. Em documentos matemáticos, a prioridade continua sendo legibilidade: número de página claro, marcas de capítulo ou seção quando úteis, e pouco ruído visual. Um bom critério é perguntar se a informação repetida ajuda o leitor a se localizar. Se não ajuda, provavelmente deve sair.

\section{Delimitadores e notação pareada}

O uso automático de \verb|\left| e \verb|\right| parece conveniente, mas nem sempre produz o melhor resultado. Esses comandos escolhem o tamanho mecanicamente e podem criar delimitadores grandes demais em expressões pequenas. Em muitos casos, tamanhos explícitos como \verb|\bigl|, \verb|\bigr|, \verb|\Bigl| e \verb|\Bigr| produzem resultado mais controlado.

Para valor absoluto e norma, \texttt{amsmath} recomenda comandos distintos para barras esquerda e direita:

\begin{lstlisting}[language={[latex]TeX}, caption={Delimitadores com amsmath}]
\newcommand{\abs}[1]{\lvert#1\rvert}
\newcommand{\norm}[1]{\lVert#1\rVert}

\[
\abs{x-y} \leq \norm{u-v}.
\]
\end{lstlisting}

O \texttt{mathtools} generaliza essa ideia com \verb|\DeclarePairedDelimiter|:

\begin{lstlisting}[language={[latex]TeX}, caption={Delimitadores com mathtools}]
\DeclarePairedDelimiter{\abs}{\lvert}{\rvert}
\DeclarePairedDelimiter{\norm}{\lVert}{\rVert}

\[
\abs{x}
\qquad
\abs*{\frac{x}{1+x^2}}
\qquad
\norm[\Big]{\sum_{k=1}^n v_k}
\]
\end{lstlisting}

A versão sem estrela usa o tamanho natural. A versão estrelada usa delimitadores automáticos. A versão com argumento opcional permite escolher explicitamente o tamanho. Essa combinação preserva uma interface semântica e ainda permite ajuste tipográfico local.

Para notações com mais estrutura, como produto interno ou conjuntos por compreensão, \texttt{mathtools} fornece \verb|\DeclarePairedDelimiterX|:

\begin{lstlisting}[language={[latex]TeX}, caption={Produto interno e conjunto por compreensão}]
\DeclarePairedDelimiterX{\inner}[2]{\langle}{\rangle}{#1,#2}

\providecommand{\given}{}
\newcommand{\SetSymbol}[1][]{%
  \nonscript\:#1\vert\allowbreak\nonscript\:\mathopen{}}
\DeclarePairedDelimiterX{\Set}[1]{\{}{\}}{%
  \renewcommand{\given}{\SetSymbol[\delimsize]}#1}

\[
\inner{u}{v}
\qquad
\Set*{x\in X \given f(x)>0}
\]
\end{lstlisting}

Esse exemplo é mais avançado, mas ilustra uma ideia central: a macro deve expressar a notação matemática, não apenas economizar digitação.

\section{Matrizes, casos e pontos}

O \texttt{amsmath} oferece ambientes de matrizes com delimitadores prontos: \texttt{pmatrix}, \texttt{bmatrix}, \texttt{vmatrix}, \texttt{Vmatrix} e outros. Para matrizes pequenas dentro de texto, \texttt{smallmatrix} evita uma construção grande demais para a linha.

\begin{lstlisting}[language={[latex]TeX}, caption={Matrizes com amsmath}]
\[
A =
\begin{pmatrix}
1 & 0\\
0 & 1
\end{pmatrix},
\qquad
\det A =
\begin{vmatrix}
a & b\\
c & d
\end{vmatrix}.
\]
\end{lstlisting}

O \texttt{mathtools} também fornece variantes úteis, como \texttt{dcases}, em que o conteúdo dos casos é composto em estilo de display:

\begin{lstlisting}[language={[latex]TeX}, caption={Casos em estilo de display}]
\[
f(x) =
\begin{dcases}
\frac{\sin x}{x}, & x\neq 0,\\
1, & x = 0.
\end{dcases}
\]
\end{lstlisting}

Para reticências, comandos semânticos como \verb|\dotsc|, \verb|\dotsb|, \verb|\dotsm| e \verb|\dotsi| comunicam o contexto: vírgulas, operadores binários, multiplicação ou integrais. Essa distinção permite que o pacote escolha a forma tipográfica apropriada.

\begin{lstlisting}[language={[latex]TeX}, caption={Reticências semânticas}]
\[
a_1, a_2, \dotsc, a_n
\qquad
a_1 + a_2 + \dotsb + a_n
\qquad
a_1 a_2 \dotsm a_n
\]
\end{lstlisting}

\section{Um exemplo de refatoração}

O trecho abaixo compila, mas mistura decisões tipográficas, notação e estrutura:

\begin{lstlisting}[language={[latex]TeX}, caption={Trecho matemático frágil}]
\begin{eqnarray}
||u+v||^2 &=& <u+v,u+v>\\
&=& ||u||^2 + 2<u,v> + ||v||^2.
\end{eqnarray}
\end{lstlisting}

Uma versão mais robusta usa ambiente adequado, macros semânticas e delimitadores declarados:

\begin{lstlisting}[language={[latex]TeX}, caption={Trecho refatorado}]
\DeclarePairedDelimiter{\norm}{\lVert}{\rVert}
\DeclarePairedDelimiterX{\inner}[2]{\langle}{\rangle}{#1,#2}

\begin{align}
\norm{u+v}^2
  &= \inner{u+v}{u+v}\\
  &= \norm{u}^2 + 2\inner{u}{v} + \norm{v}^2.
\end{align}
\end{lstlisting}

A diferença não é apenas visual. Na segunda versão, o ambiente expressa que há um cálculo alinhado; os comandos \verb|\norm| e \verb|\inner| expressam a notação; a troca de estilo dos delimitadores pode ser feita em um único ponto do preâmbulo.

\section{Critérios práticos}

Ao escrever matemática em um documento técnico, algumas perguntas ajudam a escolher a estrutura:

\begin{itemize}
    \item A fórmula é uma equação isolada ou parte de um cálculo?
    \item Cada linha precisa de número?
    \item Há um ponto matemático natural de alinhamento?
    \item A notação aparece muitas vezes no documento?
    \item O delimitador deve ter tamanho automático ou tamanho controlado?
    \item O texto dentro da fórmula é uma palavra, uma justificativa ou parte da notação?
    \item A informação repetida em cabeçalhos e rodapés ajuda a navegação ou apenas adiciona ruído?
\end{itemize}

Essas perguntas evitam uma escrita matemática puramente visual. Em \LaTeX, a qualidade do resultado depende de escolher estruturas que correspondem ao papel matemático do trecho.
